\documentclass[10pt,a4paper]{amsart}
\usepackage[margin=19mm]{geometry}
\usepackage{amsmath,amssymb,mathtools}
\usepackage[scr=rsfs,cal=euler]{mathalfa}
\usepackage{xfrac}
\usepackage[unicode,hidelinks,bookmarks=false]{hyperref}
\newcommand{\ie}{i.e.\ }
\newcommand{\cf}{cf.\ }
\newcommand{\resp}{respectively\ }
\newcommand{\Subsection}{\subsection}
\providecommand{\nobreakdash}{\nobreak\hbox{-}}
\setlength{\emergencystretch}{2em}
\multlinegap0pt
\usepackage{dsfont}
\usepackage{amssymb}
\usepackage{enumerate}
\newtheorem{proposition}{Proposition}
\newtheorem{lemma}{Lemma}
\title{Curvature characterization of Hermitian manifolds with Bismut parallel torsion}
\author{Quanting Zhao, Fangyang Zheng$^{\ast}$}
\date{Source: arXiv:2407.10497v3 (CC BY 4.0)}
\begin{document}
\maketitle

\noindent\emph{Source and licence.} Curvature characterization of Hermitian manifolds with Bismut parallel torsion. Version arXiv:2407.10497v3 (CC BY 4.0), distributed by arXiv under the Creative Commons Attribution 4.0 International licence, http://creativecommons.org/licenses/by/4.0/, as recorded in the arXiv licence metadata for this exact version and shown on the abstract page https://arxiv.org/abs/2407.10497v3. Full licence notice: \texttt{licenses/arxiv-2407.10497-CC-BY-4.0.txt}.\par\smallskip

\noindent\textit{Editorial scope.} Counted unit: the article's proposition prop1.5 (source0.tex lines 180--192). The article's complete original proof of it is delivered with it (source0.tex lines 997--1050, one \texttt{proof} environment of 54 lines, which the article places away from the statement). No mathematics is written, repaired or re-derived: the statement and the proof are the article's own lines, in the article's own notation, and no retained line is substituted or re-flowed.  Retained uncounted as the source-local context the proof needs: lines 515--532; lines 630--640; lines 658--674.  Nothing was omitted from any retained span, so the package carries no omission record.  No retained line contains a citation, so the wrapper carries no bibliography and no bibliography entry.  No retained line contains a figure environment, an image-inclusion command or drawing code, so no image is delivered and the package declares no asset dependency.  Editorial formatting only, in addition: an AMS \texttt{amsart} preamble that restates the article's own declarations and macros used by the retained lines, namely the environments \texttt{proposition}, \texttt{lemma} on separate counters, together with the standard packages those lines need.  The retained environments print in this document as Proposition 1, Lemma 1 in order of appearance, whereas the article prints the same items with the numbering of its own full document.  The complete native source of the article as distributed in the arXiv e-print for this exact version is bundled unchanged as \texttt{sources/162319/source0.tex}.
\begin{lemma} \label{lemma3.1}
Let $(M^n,g)$ be a Hermitian manifold. Then under any local unitary frame the following identities hold:
\begin{eqnarray}
&& R^{b}_{ijk\bar{\ell}} \ = \
\big( T^{\ell}_{kj,i}-T^{\ell}_{ki,j} \big) +\sum_r \big( T^r_{ij}T^{\ell}_{rk} + T^r_{jk}T^{\ell}_{ri} + T^r_{ki}T^{\ell}_{rj} \big) ,       \label{curv20}\\
 && R^b_{i\bar{j}k\bar{\ell}}-R^c_{i\bar{j}k\bar{\ell}} \ = \ \big( T^{\ell}_{ik,\bar{j}} + \overline{T^{k}_{j\ell,\bar{i}}} \big)  -  \sum_r  \big( T^{\ell}_{kr}\overline{T^{i}_{jr}} + T^{j}_{ir}\overline{T^k_{\ell r}}
+ T^r_{ik}\overline{T^r_{j\ell}} - T^{\ell}_{ir}\overline{T^{k}_{jr}} \big),     \label{curv11}\\
&& T^k_{ij,\ell} + T^k_{j\ell,i} + T^k_{\ell i,j} \ = \ 2\sum_r  \big( T^r_{\ell i}T^k_{r j} + T^r_{j\ell}T^k_{ri} + T^r_{ij}T^k_{r \ell} \big) ,      \label{pT} \\
&& T^j_{ik,\bar{\ell}} + \overline{T^i_{j\ell,\bar{k}}} - \overline{T^k_{j\ell,\bar{i}}} \ \, = \ \,   -\, \sum_r \big( T^r_{ik}\overline{T^r_{j\ell}} + T^j_{ir} \overline{T^k_{\ell r}} - T^j_{kr}\overline{T^i_{\ell r}}
-T^\ell_{ir} \overline{T^k_{jr}} + T^\ell_{kr}\overline{T^i_{jr}} \big)       \label{dbT} \\
&& \hspace{4cm} + \ \,\big( R^{b}_{i\bar{\ell}k\bar{j}}-R^b_{k\bar{\ell}i\bar{j}} \big), \notag \\
&& R^{b}_{ijk\bar{\ell},p} + R^{b}_{pik\bar{\ell},j} +R^{b}_{jpk\bar{\ell},i} \ = \
 -\,\sum_r \big( R^b_{irk\bar{\ell}}T_{jp}^r + R^b_{jrk\bar{\ell}}T_{pi}^r + R^b_{prk\bar{\ell}}T_{ij}^r \big) , \label{dR_30}  \\
&& R^{b}_{ipk \bar{\ell},\bar{q}} - R^{b}_{i\bar{q}k \bar{\ell}, p} + R^{b}_{p \bar{q} k \bar{\ell}, i} \  = \ \sum_r \{  R^b_{r\bar{q}k\bar{\ell}}T^r_{ip}  - R^b_{p\bar{r}k\bar{\ell}}T^q_{ir}   + R^b_{i\bar{r}k\bar{\ell}}T^q_{pr} +
\label{dR_21} \\
&& \hspace{4.8cm} +  \ R^b_{prk\bar{\ell}} \overline{T^i_{qr}}  - R^b_{irk\bar{\ell}}\overline{T^p_{qr}}\} ,\notag \end{eqnarray}
for any $i,j,k,\ell,p,q$, where indices after comma stand for covariant derivatives with respect to $\nabla^{b}$.
\end{lemma}

\begin{proposition}\label{Tderivative}
Let $(M^n,g)$ be a Hermitian manifold, with $\omega$ its K\"ahler form.
Then it holds that
\begin{eqnarray}
\nabla^b_{1,0} T^c=0 & \Longleftrightarrow &
R^b_{ijk\bar{\ell}}=0\quad \Longrightarrow \quad \sum_r(T^r_{\ell i}T^k_{r j} + T^r_{j\ell}T^k_{ri} + T^r_{ij}T^k_{r \ell})=0,\label{parallel_10}\\
\nabla^b_{0,1} T^c=0 & \Longleftrightarrow &
R^b_{i \bar{j} k \bar{\ell}} - R^b_{k \bar{j} i \bar{\ell}} = -P_{ik}^{j \ell} \label{parallel_01},
\end{eqnarray}
where $\nabla^b_{1,0} T^c$ and $\nabla^b_{0,1} T^c$ are respectively the $(1,0)$- and $(0,1)$-components of $\nabla^bT^c$.
\end{proposition}

\begin{proposition} \label{prop3.6}
Let $(M^n,g)$ be a BTP manifold. Then it holds that, for any $i,j,k,\ell$,
\begin{enumerate}
\item\label{R20=0} $R^b_{ijk\bar{\ell}}=0$,
\item\label{R11_12=34} $R^b_{i \bar{j} k \bar{\ell}} = R^b_{k \bar{\ell} i \bar{j}}$,
\item\label{Q_pll} $\nabla^b Q=0$,
\item\label{R=0} $R^b_{xy \chi w}=0$,
\end{enumerate}
where $x,y,w$ are any tangent vector and $\chi$ is the dual vector field of Gauduchon's torsion $1$-form $\eta$.
Also, under the BTP condition, we have
\begin{eqnarray}
&& \partial \overline{\partial} \omega =0 \quad \Longleftrightarrow \quad   P_{ik}^{j \ell}=0  \quad \Longleftrightarrow \quad
R^b_{i \bar{j} k \bar{\ell}} = R^b_{k \bar{j} i \bar{\ell}}, \label{eq_plcld}\\
&& \partial \eta=0,\quad \overline{\partial} \eta = -\sum_{i,j,k}\eta_k \overline{T^i_{jk}} \varphi_i \wedge \overline{\varphi}_j,
\end{eqnarray}
In particular, any BTP metric $g$ is always Gauduchon, namely, satisfies $\partial \overline{\partial} \omega^{n-1}=0$.
\end{proposition}

\begin{proposition} \label{prop1.5}
Let $(M^n,g)$ be a BTP manifold. Then under any local unitary frame $e$, the following hold:
\begin{eqnarray}
&&  \sum_r \big( T^r_{ij}T^{\ell}_{r k} + T^r_{jk}T^{\ell}_{ri} + T^r_{ki}T^{\ell}_{rj}\big) =0,  \ \ \ \ \ \forall \ 1\leq i,j,k,\ell \leq n, \label{eq:1.5} \\
&& 2 \sum_{r,s} T^r_{si} A_{r\bar{s}} =  \sum_{r,s} T^r_{si} B_{r\bar{s}} , \ \ \ \ \ \forall \ 1\leq i\leq n, \label{eq:1.6} \\
&& Q_{i\bar{j}k\bar{\ell}}  = \sum_r \big( T^j_{kr} \overline{ T^i_{\ell r} } + T^{\ell}_{ir} \overline{ T^k_{j r} } - T^j_{ir} \overline{ T^k_{\ell r} } - T^{\ell}_{kr} \overline{ T^i_{j r} } - T^r_{ik} \overline{ T^r_{j\ell } } \big) ,\ \ \ \ \ \forall \ 1\leq i,j,k,\ell \leq n, \label{eq:1.7} \\
&& \partial \eta =0, \ \ \ \overline{\partial } \eta = - \sum_{i,j} \overline{\phi^i_j }\,\varphi_i \wedge \overline{\varphi}_j, \label{eq:1.8}\\
&& R^b_{i\bar{j}k\bar{\ell}} - R^c_{i\bar{j}k\bar{\ell}} = \sum_r \big( T^{\ell}_{ir} \overline{T^{k}_{jr}} -T^{r}_{ik} \overline{T^{r}_{j\ell }} - T^{j}_{ir} \overline{T^{k}_{\ell r}} - T^{\ell}_{kr} \overline{T^{i}_{jr}}\big) , \ \ \ \ \forall \ 1\leq i,j,k,\ell \leq n, \label{eq:1.9}\\
&& \sum_r \big( \phi^r_i T^j_{rk} + \phi^r_k T^j_{ir} - \phi^j_r T^r_{ik} \big) = 0, \ \  \ \ \forall \ 1\leq i,j,k\leq n, \label{eq:1.10}\\
&& A, \,B, \,\phi, \,\phi^{\ast} \, \mbox{commute with each other}. \label{eq:1.11}
\end{eqnarray}
Here $\phi^{\ast}$ stands for $ (\phi^{\ast})^j_i = \overline{\phi^i_j}$, $\varphi$ is the coframe dual to $e$, while $R^b$ and $R^c$ denote the curvature tensor of the Bismut and Chern connection, respectively.
\end{proposition}

\begin{proof}[{\bf Proof of Proposition \ref{prop1.5}.}]
Let $(M^n,g)$ be a BTP manifold. Note that the equalities (\ref{eq:1.5}) and (\ref{eq:1.7})\,--\,(\ref{eq:1.9})
have already been established in Lemma \ref{lemma3.1}, Propositions \ref{Tderivative} and \ref{prop3.6}.
Multiple the equality (\ref{eq:1.5}) by $\overline{T_{jk}^\ell}$ and sum up $j,k,\ell$, then the equality (\ref{eq:1.6}) is also established.
So we will start with the proof of (\ref{eq:1.10}).

Fix any $p\in M$. Choose a local unitary frame $e$ in a neighborhood of $p$
so that the matrix of Bismut connection vanishes at $p$, i.e. $\theta^b|_p=0$.
Denote by $\varphi$ the coframe dual to $e$. Then the Chern connection matrix satisfies $\theta|_p=-\gamma|_p$.
Hence by the structure equation
$ d\varphi = - \,^t\!\theta \wedge \varphi +\tau$, we have
$$ \partial \varphi_i = -\frac{1}{2}\sum_{r,s}T^i_{rs}\varphi_r \wedge \varphi_s, \ \ \ \overline{\partial}\varphi_j = - \sum_{t,r} \overline{T^t_{jr}} \, \overline{\varphi}_r \wedge \varphi_t, \ \ \ \ \ \mbox{at} \ \ p. $$
The conjugation of the second equality of (\ref{eq:1.8}) implies that $\partial \overline{\eta} = \sum_{i,j} \phi^j_i \varphi_i \wedge \overline{\varphi}_j$. After taking $\partial $ again and using the fact that $\nabla^b\phi =0$, we know that at $p$
\begin{eqnarray*}
 0 & = &  \partial^2\overline{\eta} = \sum_{i,j} \phi^j_i \big\{ -\frac{1}{2}\sum_{r,s} T^i_{rs} \varphi_r\varphi_s \wedge \overline{\varphi}_j - \varphi_i \wedge  \sum_{t,r} T^t_{jr} \varphi_r \overline{\varphi}_t \big\} \\
 & = & \frac{1}{2} \sum_{i,j,k,r} \big( - \phi^j_r T^r_{ik} + \phi^r_i T^j_{rk} + \phi^r_k T^j_{ir} \big) \,\varphi_i\varphi_k\overline{\varphi}_j .
\end{eqnarray*}
Thus (\ref{eq:1.10}) is established. Next we will prove (\ref{eq:1.11}).
Let us multiply $\overline{\eta}_k$ onto both sides of the equality (\ref{eq:1.7}) and sum up $k$.
The equality (4) of Proposition \ref{prop3.6} indicates the left hand side becomes zero, so we end up with
\begin{equation} \label{eq:4.1}
 0 = \sum_r \{ -\phi^j_r \overline{T^i_{\ell r}} + \phi^{\ell }_r \overline{T^i_{j r}}  - \phi^r_i \overline{T^r_{j\ell }}  \} , \ \ \ \ \ \forall \ i,j,\ell.
 \end{equation}
Multiply by $\eta_{\ell}$ onto the above equation and sum up $\ell$.
By utilizing the fact that $\sum_{\ell} T^{\ell}_{rs}\eta_{\ell} =0$ which can be derived from (\ref{eq:1.5}), we get
$$ 0 = \sum_r \{ \phi^j_r \overline{\phi^i_r } - \phi^r_i \overline{\phi^r_j } \} .$$
that is, $\phi$ commutes with $\phi^{\ast}$. If we multiply $T^s_{j\ell }$ onto (\ref{eq:4.1}) and sum up $j$, $\ell$,
then we have \[\sum_r \phi_i^r B_{r\bar{s}} = 2 \sum_{j,\ell,r} \phi^j_r T^s_{j \ell} \overline{T^i_{r \ell}}
=\sum_{r} B_{i\bar{r}} \phi_r^s.\]
Similarly, if we multiply $T^i_{\ell s}$ onto (\ref{eq:4.1}) and sum up $i$, $\ell$, then we have
\[\sum_r A_{s\bar{r}}\phi_r^j
=\sum_{i,\ell,r}(\phi^{\ell}_r T^i_{\ell s} \overline{T^i_{jr}} + \phi_i^r T^i_{s \ell}\overline{T^r_{j \ell}})
=\sum_r \phi_s^r A_{r \bar{j}}.\]
Note that here we have used (\ref{eq:1.5}) in obtaining  both of the above equalities. Hence, $\phi$ commutes with both $A$ and $B$.

Finally we will show that $A$ commutes with $B$.
To this end, let us choose our local unitary frame $e$ so that $A$ is diagonal, which can be decomposed into
several blocks, where diagonal entries remain the same in each block, while different blocks enjoy distinct representing entries.
Since $\nabla^b A=0$, it yields that
$$ d A_{i\bar{j}} = \sum_r \{  A_{r\bar{j}} \theta^b_{ir} - A_{i\bar{r}} \theta^b_{rj} \} , \ \ \ \ \ \ \forall \  i,j. $$
This implies that $\theta^b_{ij}=0$ whenever $A_{i\bar{i}} \neq A_{j\bar{j}}$,
so $\theta^b$ is block-diagonal with respect to those blocks of $A$.
By the structure equation $\Theta^b=d\theta^b - \theta^b\wedge \theta^b$,
the matrix of Bismut curvature $\Theta^b$ is also block-diagonal.
This means that $R^b_{i\bar{j}\ast \bar{\ast}} = R^b_{\ast \bar{\ast}i\bar{j}} =0$
whenever $A_{i\bar{i}} \neq A_{j\bar{j}}$, i.e. $i$ and $j$ are from different blocks.
In the meantime, for any $i$ and $j$ from different blocks,
$$ \mathrm{Ric}(Q)_{i \bar{j}} = \sum_r \big\{ R^b_{i\bar{j}r\bar{r}} - R^b_{r\bar{j}i\bar{r}} \big\} = - \sum_r R^b_{r\bar{j}i\bar{r}} =0,$$
since $r$ cannot belong to the same block with $i$ and $j$ simultaneously.
This shows that $\mathrm{Ric}(Q)$ is block-diagonal with respect to those $A$-blocks, hence it commutes with $A$.
On the other hand, by letting $k=\ell$ and summing up in (\ref{eq:1.7}), we get $\mathrm{Ric}(Q)=B-\phi -\phi^{\ast}$.
Since we already showed that $A$ commutes with $\phi$ hence $\phi^{\ast}$, we know that $A$ commutes with $B$ as well.
This completes the proof of Proposition \ref{prop1.5}.
\end{proof}

\end{document}
